\chapter{Parcours d'intégration au sein de l'entreprise}\label{ch:parcours-d'integration-au-sein-de-l'entreprise}


\section{Entretien réglementaire avec les différents interlocuteurs}\label{sec:entretien-reglementaire-avec-les-differents-interlocuteurs}

    La première partie de l'alternance a été marquée par un certain nombre d'entretiens réglementaires et d'évènements de différentes natures.

    \begin{itemize}
        \item À l'échelle du \gls{CEA} de Grenoble :
        \begin{itemize}
            \item Une journée d'accueil des alternants organisée par le service du personnel, permettant de découvrir les activités du Centre ainsi que ses nombreux équipements
            \item Une visite médicale
            \item Une journée d'accueil dédiée à la sécurité \textit{(à venir car incompatible avec le calendrier d'alternance jusqu'alors)}
        \end{itemize}
        \item À l'échelle de l'\gls{INES}, avec une série de présentations par différents interlocuteurs :
        \begin{itemize}
            \item Assistante de laboratoire
            \item Responsable sécurité
            \item Responsable qualité
            \item Responsable de sécurité des systèmes d'information
            \item Correspondant scientifique
            \item Chef d'installation
        \end{itemize}
    \end{itemize}
\section{Accueil par l'équipe et immersion progressive}\label{sec:accueil-par-l'equipe-et-immersion-progressive}

    L'ensemble des procédures administratives a été complété par un accueil chaleureux de l'ensemble du \gls{LELA}, notamment par une visite
de l'\gls{INES} et une présentation à l'équipe par l'encadrant de l'alternance en amont de celle-ci.
La découverte du \gls{CEA} par l'alternant s'est faite aussi et surtout par de nombreuses discussions informelles autour des différents projets menés par le laboratoire, des travaux de recherche des collègues et de leur quotidien.
En outre, plusieurs réunions techniques avec quelques membres du laboratoire ont permis la compréhension des enjeux de l'alternance et leurs liens avec d'autres activités, tout en favorisant la cohésion d'équipe.

Enfin, le travail de cohésion entre les membres de l'\gls{INES} est appuyé par son association Act'\gls{INES} proposant à ses adhérents de nombreuses activités extraprofessionnelles de natures sportive et culturelle.
